\chapter{General Experiments and Blackwell's Martingale Construction}
\label{sec:p17-general-blackwell-comparison}

An observation need not have a finite range. For finitely many states,
its relevant likelihood information can still be represented by a
probability distribution on a finite-dimensional simplex. Blackwell's
general construction compares such distributions through a joint
probability law whose conditional mean is prescribed.

The presentation below supplies the two approximation chains, their
compatible stochastic transformations, the countable joint construction,
both martingale limits and the final probability kernel in the source's
Theorem~8. It also supplies the standard-measure and decision comparisons
of Theorems~1--2 with explicit measurable-space hypotheses
\cite[pp.~265--269]{Blackwell1953EquivalentExperiments}.

\section{Convex tests and probability kernels}
\label{sec:p17-general-convex-tests}

\begin{hbdefinition}{def:p17-general-convex-order}{Bounded laws and their comparison}
Let \(d\geq1\), and let \(\mu,\nu\) be Borel probability measures
concentrated on a nonempty bounded Borel set \(X\subseteq\mathbb R^d\).
Use the Euclidean norm \(\|x\|_2=(\sum_i x_i^2)^{1/2}\);
\(\delta_x\) denotes the point mass at \(x\), and
\(\operatorname{conv}(X)\) its set of finite convex combinations.
Write
\[
 \mu\preceq_{\mathrm{cx}}\nu
 \quad\Longleftrightarrow\quad
 \int\phi\,d\mu\leq\int\phi\,d\nu
       \quad\text{for every continuous convex }
                         \phi:\mathbb R^d\to\mathbb R.
\]
These integrals are finite because the measures have bounded support
and each test is continuous on a surrounding compact box.
The direction means that \(\nu\) is at least as dispersed under convex
tests. Affine functions and their negatives are included.

A probability kernel \(K:X\times\mathcal B(X)\to[0,1]\) is
\emph{mean preserving} if
\[
                     \int_X y\,K(x,dy)=x\qquad(x\in X),
\]
coordinate by coordinate. Its action on a measure is
\((\mu K)(A)=\int_XK(x,A)\,d\mu(x)\).
The equality is required at every input point, including points
having zero \(\mu\)-probability.
For kernels, composition is in row order:
\((KL)(x,A)=\int L(y,A)K(x,dy)\).
Finite kernels are row-stochastic matrices with the same convention.
The mean-preserving definition also applies to kernels between two
different bounded Borel subsets of the same Euclidean space.
\index{convex order}\index{mean-preserving kernel}
\end{hbdefinition}

\begin{hblemma}{thm:p17-compact-convex-tools}{proved=annex:p17.compact-convex-tools}{Affine approximation and Jensen on compact convex sets}
Let \(C\) be a nonempty compact convex subset of a finite-dimensional
real space. Every continuous convex \(\phi:C\to\mathbb R\) is uniformly
approximated from below by maxima of finitely many affine functions
defined on the whole ambient space.

Every Borel probability measure \(\lambda\) on \(C\) has its barycenter
in \(C\) and satisfies
\[
                         \phi\left(\int_Cx\,d\lambda(x)\right)
                                   \leq\int_C\phi(x)\,d\lambda(x).
\]
Consequently, for finitely supported probabilities, the global convex
tests in the preceding definition suffice to apply
\cref{thm:p17-finite-convex-transport}, including its zero-mass rows.
\end{hblemma}

\begin{hblemma}{thm:p17-continuous-law-identification}{proved=annex:p17.continuous-law-identification}{Identifying a law by continuous tests}
Two Borel probabilities on \(\mathbb R^d\) with equal integrals of
every bounded continuous function are equal.
If all the measures in question are concentrated on one compact set,
it suffices to use continuous functions on that set.
For such functions, their uniform modulus of continuity controls any
approximation that moves a point by at most a specified distance.
\end{hblemma}

\section{The two finite approximation chains}
\label{sec:p17-blackwell-approximations}

\begin{hbdefinition}{def:p17-dyadic-transport-data}{Cells, centroids and normalized vertex weights}
For \(N\geq1\), put \(h_N=2^{-N}\). For
\(a=(a_1,\ldots,a_d)\in\mathbb Z^d\), define
\[
                 C_{N,a}=\prod_{i=1}^d[h_Na_i,h_N(a_i+1)).
\]
When \(\mu(C_{N,a})>0\), let
\[
      c_{N,a}=\frac{\int_{C_{N,a}}t\,d\mu(t)}{\mu(C_{N,a})},
      \qquad
      \mu_N=\sum_{a:\,\mu(C_{N,a})>0}
                                      \mu(C_{N,a})\,\delta_{c_{N,a}}.
\]
Combine coincident centroids when writing this as a probability on its
positive-mass support \(A_N\). No centroid is assigned to a zero-mass cell.

For \(t\in C_{N,a}\), put
\[
 \lambda_{N,i}(t)=\frac{t_i-h_Na_i}{h_N},\qquad
 v_{N,a,\epsilon}=h_N(a+\epsilon),\qquad
 \epsilon\in\{0,1\}^d,
\]
and define
\[
 \begin{split}
 w_{N,\epsilon}(t)
   &=\prod_{i=1}^d
       \begin{cases}
          1-\lambda_{N,i}(t),&\epsilon_i=0,\\
          \lambda_{N,i}(t),&\epsilon_i=1,
        \end{cases}\\
 U_N(t,\cdot)&=\sum_{\epsilon\in\{0,1\}^d}
                         w_{N,\epsilon}(t)\,\delta_{v_{N,a,\epsilon}},
 \qquad \nu_N=\nu U_N.
 \end{split}
\]
Let \(B_N\) be the positive-mass support of \(\nu_N\).
The factor \(h_N^{-1}\) is part of each barycentric coordinate.
Cell membership is unique, including at a dyadic grid point.
\end{hbdefinition}

\begin{hbtheorem}{thm:p17-centroid-refinement}{proved=annex:p17.centroid-refinement}{The increasing centroid chain}
The \(\mu_N\) are finitely supported probabilities in one fixed compact
box. There are mean-preserving finite kernels
\(T_N:A_N\to A_{N+1}\) with
\[
                 \mu_NT_N=\mu_{N+1},\qquad
                 \mu_N\preceq_{\mathrm{cx}}\mu_{N+1}
                         \preceq_{\mathrm{cx}}\mu.
\]
For every continuous \(f\) on a surrounding compact box \(K\),
\[
        \left|\int f\,d\mu_N-\int f\,d\mu\right|
                            \leq\omega_{f,K}(\sqrt d\,h_N),
\]
where
\(\omega_{f,K}(r)=\sup\{|f(s)-f(t)|:s,t\in K,\ \|s-t\|_2\leq r\}\).
Coincident centroids are treated by the conditional law of their
actual coarse/fine pair, rather than by selecting a single parent cell.
\end{hbtheorem}

\begin{hbtheorem}{thm:p17-vertex-refinement}{proved=annex:p17.vertex-refinement}{The decreasing vertex chain}
Every \(U_N\) is a Borel probability kernel on \(\mathbb R^d\), with
barycenter \(t\) at every \(t\). Its refinements satisfy
\[
                      U_{N+1}U_N=U_N.
\]
For \(N\geq2\), the restriction \(V_N=U_{N-1}|_{B_N}\) takes values
in \(B_{N-1}\), is mean preserving, and obeys
\[
                \nu_NV_N=\nu_{N-1}.
\]
For every \(N\geq1\),
\(\nu\preceq_{\mathrm{cx}}\nu_{N+1}\preceq_{\mathrm{cx}}\nu_N\).
All the \(\nu_N\) lie in one compact box \(K\). For every continuous
\(f:K\to\mathbb R\),
\[
        \left|\int f\,d\nu_N-\int f\,d\nu\right|
                            \leq\omega_{f,K}(\sqrt d\,h_N).
\]
These are convergence statements for continuous tests.
\end{hbtheorem}

\begin{hbtheorem}{thm:p17-compatible-transports}{proved=annex:p17.compatible-transports}{A simultaneous compatible family of finite transformations}
Assume \(\mu\preceq_{\mathrm{cx}}\nu\). There exist mean-preserving
finite kernels \(W_N:A_N\to B_N\) such that
\[
                 \mu_NW_N=\nu_N,\qquad
                 W_N=T_NW_{N+1}V_{N+1}\quad(N\geq1).
\]
The proof first chooses arbitrary finite comparisons at distant levels,
transports them to each fixed level, and extracts a single diagonal
subsequence. The equation is a matrix identity at every row.
It is Blackwell's compatibility condition written in the declared
row-composition convention.
\end{hbtheorem}

\section{The joint process and the general transport}
\label{sec:p17-blackwell-joint-process}

\begin{hbtheorem}{thm:p17-blackwell-joint-process}{proved=annex:p17.blackwell-joint-process}{The complete two-sided construction and its limits}
Assume \(\mu\preceq_{\mathrm{cx}}\nu\), and choose the compatible
kernels \(T_N,V_N,W_N\) just constructed.
There is a probability space carrying bounded random vectors
\((L_N,R_N)_{N\geq1}\) such that, for every \(N\), the ordered string
\[
                    L_1,\ldots,L_N,R_N,\ldots,R_1
\]
has initial law \(\mu_1\) and successive kernels
\(T_1,\ldots,T_{N-1},W_N,V_N,\ldots,V_2\).
Its marginals are \(\mu_N,\nu_N\).

The \(L_N\) form a forward martingale for
\(\mathcal F_N=\sigma(L_1,\ldots,L_N)\).
The \(R_N\) form a backward martingale for
\(\mathcal G_N=\sigma(R_N,R_{N+1},\ldots)\).
Their almost-sure and \(L^1\) limits \(L_*,R_*\) have laws \(\mu,\nu\)
and satisfy
\[
                          E(R_*\mid L_*)=L_*.
\]
The assertion uses the full countable extension and both conditional-limit
theorems; all relevant filtrations and consistency equations are supplied
in the annex.
\end{hbtheorem}

\begin{hbtheorem}{thm:p17-general-convex-transport}{proved=annex:p17.general-convex-transport}{Blackwell's general bounded transport theorem}
For the bounded Borel set \(X\) and laws \(\mu,\nu\) above,
\(\mu\preceq_{\mathrm{cx}}\nu\) if and only if there is a
mean-preserving Borel probability kernel \(K:X\to X\) with
\[
                                  \mu K=\nu.
\]
Equivalently there is a joint law of \(L,R\), with these respective
marginals, such that \(E(R\mid L)=L\).

In particular the source's premise of domination for every continuous
convex function on \(\operatorname{conv}(X)\) implies this kernel
conclusion: restricting a global convex test gives one of those tests.
Thus the full bounded-set implication of source Theorem~8 is supplied.
The kernel is repaired measurably on its exceptional input set by
\(\delta_x\), giving the barycenter identity at every \(x\in X\)
\cite[pp.~268--269]{Blackwell1953EquivalentExperiments}.
\end{hbtheorem}

\begin{hbtheorem}{thm:p17-transport-calculations}{proved=annex:p17.transport-calculations}{Weights, open-set limits and an explicit continuous transport}
\begin{enumerate}
\item At level \(N=1\) in two dimensions, for
      \(t=(1/4,3/8)\), the four vertices
      \((0,0),(1/2,0),(0,1/2),(1/2,1/2)\) have weights
      \(1/8,1/8,3/8,3/8\), respectively.
\item Let \(\mu=\frac12\delta_0+\frac12\operatorname{Unif}(0,1)\).
      Its centroid approximations have
      \(\mu_N((0,\infty))=1\) for every \(N\), whereas
      \(\mu((0,\infty))=1/2\).
\item For \(X=[0,1]\), uniform \(\mu\), and
      \(\nu=\frac12\delta_0+\frac12\delta_1\), the kernel
      \[
                K(x,\cdot)=(1-x)\delta_0+x\delta_1
      \]
      is a mean-preserving transformation from \(\mu\) to \(\nu\).
      The reverse comparison fails, as the convex test \(x^2\) shows.
\item For \(\mu=\nu=\delta_0\) on \(X=[0,1]\), an arbitrary conditional
      kernel may use \(\delta_0\) at every input and still represent the
      joint law. That version is not mean preserving at \(x>0\).
      The version \(K(x,\cdot)=\delta_x\) meets the every-input condition.
\end{enumerate}
\end{hbtheorem}

\section{General observations and standard measures}
\label{sec:p17-general-experiments}

\begin{hbdefinition}{def:p17-measurable-experiment}{A finite-state experiment on a measurable observation space}
Let \(S=\{1,\ldots,n\}\), \(n\geq1\).
An experiment \(\mathsf E=(P_s)_{s\in S}\) on a nonempty measurable
space \((Z,\mathcal Z)\) is an ordered family of probabilities.
Fix a prior \(\pi_s>0\), \(\sum_s\pi_s=1\), and put
\[
      \rho=\sum_s\pi_sP_s,\qquad
      p_s(z)=\pi_s\frac{dP_s}{d\rho}(z).
\]
Choose the finite measurable versions so that
\(p(z)\in\Delta(S)\) everywhere, using a fixed point of the simplex
on a common \(\rho\)-null set. Its law under \(\rho\) is
\(\mu_{\mathsf E,\pi}\).

For a finite action set \(D\ne\varnothing\) and real payoffs \(u(d,s)\),
a randomized policy is a measurable family
\(\delta(d\mid z)\geq0\) summing to one. Define
\[
 \begin{split}
 J(\pi,\mathsf E,u,\delta)
     &=\sum_s\pi_s\int_Z\sum_{d\in D}\delta(d\mid z)u(d,s)\,dP_s(z),\\
 V(\pi,\mathsf E,u)&=\sup_\delta J(\pi,\mathsf E,u,\delta).
 \end{split}
\]
The definitions of \(J\) and \(V\) also apply to priors with zero
entries; full support is required for the posterior reconstruction above.
For a nonempty compact convex \(C\subseteq\mathbb R^n\), let
\[
 B(\mathsf E,C)=
 \left\{\left(\int f_s\,dP_s\right)_{s=1}^n:
                             f:Z\to C\text{ measurable}\right\}.
\]
These are statewise losses, with randomization already represented by
convexity of \(C\).
\index{experiment!measurable observations}\index{standard measure!experiment}
\end{hbdefinition}

\begin{hbtheorem}{thm:p17-general-standard-measure}{proved=annex:p17.general-standard-measure}{Posterior representation and implementation}
The preceding density versions exist and
\[
 \int_{\Delta(S)}q_s\,d\mu_{\mathsf E,\pi}(q)=\pi_s,\qquad
 V(\pi,\mathsf E,u)=
                \int_{\Delta(S)}\max_{d\in D}\sum_sq_su(d,s)
                                             \,d\mu_{\mathsf E,\pi}(q).
\]
An optimal measurable deterministic policy exists.

Every Borel probability \(\mu\) on \(\Delta(S)\) with barycenter \(\pi\)
is implemented by the canonical experiment
\[
                       P_s^\mu(dq)=\frac{q_s}{\pi_s}\,\mu(dq).
\]
The original and canonical experiments have exactly the same attainable
loss sets \(B(\mathsf E,C)\) for every compact convex \(C\).
No standard Borel hypothesis on the original space is required for this
loss-set statement.

For the uniform prior, \(p_s=dP_s/d(\sum_jP_j)\), and the law of \(p\)
under \((\sum_jP_j)/n\) is precisely the source's standard measure.
Its barycenter is \((1/n,\ldots,1/n)\).
\end{hbtheorem}

\begin{hblemma}{thm:p17-compact-measure-selection}{proved=annex:p17.compact-measure-selection}{Selecting a limiting law on a compact set}
For a nonempty compact \(K\subseteq\mathbb R^m\), every sequence of
Borel probabilities on \(K\) has a subsequence whose integrals of
all continuous functions converge to those of one Borel probability
on \(K\). The proof constructs consistent dyadic cell masses and
their nested-cell limit.
\end{hblemma}

\begin{hbtheorem}{thm:p17-general-risk-compactness}{proved=annex:p17.general-risk-compactness}{Closed bounded convex attainable loss sets}
For every measurable finite-state experiment and nonempty compact
convex \(C\), the set \(B(\mathsf E,C)\) is nonempty, compact and convex.
Its closedness follows from a limiting joint law of posterior and
action, followed by conditional averaging into \(C\).
\end{hbtheorem}

\begin{hbtheorem}{thm:p17-general-decision-comparison}{proved=annex:p17.general-decision-comparison}{General loss-set and posterior comparisons}
For experiments \(\mathsf E,\mathsf F\) with the same finite state set,
the following conditions are equivalent:
\begin{enumerate}
\item \(B(\mathsf E,C)\supseteq B(\mathsf F,C)\) for every nonempty
      compact convex loss set \(C\);
\item \(\mu_{\mathsf F,\pi}\preceq_{\mathrm{cx}}\mu_{\mathsf E,\pi}\)
      for one fixed full-support prior \(\pi\);
\item \(V(\pi,\mathsf E,u)\geq V(\pi,\mathsf F,u)\) for every finite
      action/payoff problem at that fixed prior;
\item the value inequality holds for every prior and every finite
      action/payoff problem.
\end{enumerate}
In the fourth condition priors with zero entries are permitted.
The equivalent loss-set condition implies the convex-order condition
at every full-support prior.
Two experiments have the same standard measure exactly when their
attainable loss sets agree for every such \(C\).
\end{hbtheorem}

\begin{hbtheorem}{thm:p17-general-garbling}{proved=annex:p17.general-garbling}{A state-independent simulation kernel}
Let \(\mathsf E=(P_s)\) be on \((Z,\mathcal Z)\) and
\(\mathsf F=(Q_s)\) on \((Y,\mathcal Y)\), with both observation
spaces standard Borel.
The conditions of \cref{thm:p17-general-decision-comparison} are also
equivalent to existence of a probability kernel \(G:Z\to Y\) with
\[
                       P_sG=Q_s\qquad(s\in S),
\]
where \(\mathsf F=(Q_s)\) is the experiment on \(Y\).
Equivalently, there is a mean-preserving kernel taking the less
informative posterior law to the more informative posterior law.
The posterior transformation runs in the opposite direction to the
observation simulation.

The proof recovers observations from their posterior values through
regular conditional laws. More generally it applies whenever the
needed observation-recovery kernel exists. The preceding loss-set
theorem itself remains valid on arbitrary measurable observation spaces.
\end{hbtheorem}

\begin{hbtheorem}{thm:p17-continuous-experiment-example}{proved=annex:p17.continuous-experiment-example}{An exact continuous-signal comparison}
For two states let \(P_1,P_2\) on \([0,1]\) have densities
\(2x\) and \(2(1-x)\), respectively. At the uniform prior the
observation law is uniform, the posterior is \((x,1-x)\), and the
optimal probability of correctly guessing the state is \(3/4\).
The single-signal experiment has value \(1/2\); it is a garbling
of this continuous experiment. The perfectly revealing two-signal
experiment garbles to it by using \(P_s\) in state-signal \(s\).

There is no simulation of this continuous experiment from the
single-signal experiment. A prior assigning probability zero to one
state cannot establish the converse comparison from value inequalities
alone.
\end{hbtheorem}

The construction above concerns bounded vector laws and finitely many
states. Unbounded transport, infinitely many states, sequential choices
and adaptive information acquisition require their additional hypotheses
and constructions. The source's separate 1951 presentation retains its
own historical identity; the proofs here reconstruct the 1953 statements.
